\chapter{Całka Wienera dla deterministycznych funkcji}
\label{ch:wiener-integral}
\paragraph{Odsyłacze międzytomowe.} Własności procesu Wienera przypomina tom II, rozdział~\ref{II-ch:wiener}, a całkowanie Monte Carlo i redukcję wariancji omawia tom IV, rozdział~\ref{IV-ch:monte-carlo}.

\section{Dlaczego nie zaczynamy od Riemanna--Stieltjesa?}
Dla gładkiej funkcji $g$ klasyczne wyrażenie
$\int_0^T f(t)\,\dd g(t)$ można rozumieć jako
$\int_0^T f(t)g'(t)\,\dd t$. Jeżeli $g$ ma jedynie
skończoną wariację, całka Riemanna--Stieltjesa nadal istnieje
dla każdej ciągłej funkcji $f$. Tymczasem dla prawie każdej
trajektorii $w(t)=W_t(\omega)$ żadne z tych założeń nie zachodzi.

Nie oznacza to, że nie da się obliczyć żadnej całki względem $w$:
na przykład przy wystarczającej regularności funkcji $f$ można
zastosować całkowanie przez części. Problem polega na tym,
że nie mamy ogólnej, ciągłej konstrukcji wzdłuż trajektorii,
która działałaby dla dowolnego $f\in L^2([0,T])$.
Co ciekawe, \emph{losowość} pozwoli ten brak regularności pokonać.

\begin{przyklad}[Dlaczego punkt próbkowania ma znaczenie]
Weźmy podział $\pi=\{0=t_0<\cdots<t_n=T\}$ i oznaczmy
$\Delta_kW=W_{t_{k+1}}-W_{t_k}$. Z tożsamości
$W_{t_{k+1}}^2-W_{t_k}^2=2W_{t_k}\Delta_kW+(\Delta_kW)^2$
wynika dokładnie
\[
\sum_{k=0}^{n-1}W_{t_k}\Delta_kW
=\frac12W_T^2-\frac12\sum_{k=0}^{n-1}(\Delta_kW)^2.
\]
W tomie II pokazaliśmy, że ostatnia suma zbiega do $T$ w $L^2$
dla podziałów deterministycznych o średnicy dążącej do zera.
Dlatego sumy \emph{lewostronne} mają granicę
$\frac12(W_T^2-T)$ w $L^2$. Gdy w każdej sumie użyjemy
punktu środkowego w \emph{wartościach procesu}, czyli
$\frac12(W_{t_k}+W_{t_{k+1}})$, otrzymamy zamiast tego
$\frac12W_T^2$ dla każdego podziału: odpowiednia suma
jest teleskopowa.

Dwie pozornie podobne reguły próbkowania dają różne odpowiedzi.
Nie jest to usterka rachunku, ale sygnał, że przy całkowaniu względem
procesu Wienera należy \emph{ustalić konwencję}. W całce Itô
stosujemy wartości funkcji podcałkowej dostępne \emph{przed} przyrostem.
\end{przyklad}

\begin{center}
\begin{tikzpicture}[>=Latex,x=1.05cm,y=0.9cm]
\draw[->] (0,0)--(8.3,0) node[right] {$t$};
\draw[->] (0,-1.45)--(0,2.35) node[above] {$w(t)$};
\draw[thick,blue!75!black] (0,0)--(.5,.3)--(1,.15)--(1.5,1.1)--(2,.5)--(2.5,.7)--(3,-.55)--(3.5,-.15)--(4.1,-.8)--(4.6,-.15)--(5.2,.75)--(5.7,.4)--(6.2,1.65)--(6.8,1.1)--(7.4,1.75);
\foreach \x in {0,1.5,3,4.6,6.2,7.4}{\draw[dashed,gray!65] (\x,-1.15)--(\x,2);}
\node[below,scale=.78] at (0,0) {$t_0$};
\node[below,scale=.78] at (1.5,0) {$t_1$};
\node[below,scale=.78] at (3,0) {$t_2$};
\node[below,scale=.78] at (4.6,0) {$t_3$};
\node[below,scale=.78] at (6.2,0) {$t_4$};
\node[below,scale=.78] at (7.4,0) {$t_5$};
\node[align=left,scale=.85] at (4.1,2.05) {Przyrost jest znany dopiero\\po osiągnięciu prawego końca};
\end{tikzpicture}
\par\smallskip\small Schemat przyrostów: reguła lewego końca nie korzysta z przyszłych wartości procesu.
\end{center}

\section{Funkcje schodkowe i pierwszy wzór całkowy}
\begin{definicja}[Deterministyczna funkcja schodkowa]
Funkcją schodkową na $[0,T]$ nazywamy funkcję postaci
$f=\sum_{k=0}^{n-1}a_k\ind{(t_k,t_{k+1}]}$, gdzie
$0=t_0<t_1<\cdots<t_n=T$ i $a_k\in\R$.
Jej wartość w punkcie $0$ nie ma znaczenia dla całek względem
miary Lebesgue'a ani dla konstruowanej całki Wienera.
Przestrzeń takich funkcji oznaczamy przez $\calS_T$.
\end{definicja}

\begin{definicja}[Całka Wienera dla funkcji schodkowej]
Dla $f=\sum_{k=0}^{n-1}a_k\ind{(t_k,t_{k+1}]}\in\calS_T$
definiujemy
\[
I_W(f):=\sum_{k=0}^{n-1}a_k(W_{t_{k+1}}-W_{t_k}).
\tag{1.1}\label{eq:wiener-simple}
\]
Wynikiem jest zmienna losowa z $L^2(\Omega)$.
\end{definicja}

\begin{lemat}[Niezależność reprezentacji]
Dla każdej funkcji $f\in\calS_T$ wartość $I_W(f)$
nie zależy od wyboru podziału ani od zapisu $f$ jako sumy
funkcji wskaźnikowych przedziałów. Odwzorowanie
$I_W:\calS_T\to L^2(\Omega)$ jest liniowe.
\end{lemat}
\begin{proof}
Jeżeli podzielimy $(t_k,t_{k+1}]$ w punkcie $u$, to
$W_{t_{k+1}}-W_{t_k}=(W_u-W_{t_k})+(W_{t_{k+1}}-W_u)$.
Wkład przedziału do sumy nie zmienia się. Dwa przedstawienia
funkcji schodkowej mają wspólne rozdrobnienie podziałów.
Po przejściu do niego współczynniki na każdym przedziale są
równe (poza punktami końcowymi), ponieważ oba przedstawienia
określają tę samą funkcję prawie wszędzie. Stąd równość sum.
Liniowość uzyskujemy przez sprowadzenie dwóch funkcji do
wspólnego podziału i dodanie współczynników.
\end{proof}

\begin{twierdzenie}[Izometria Wienera i prawo Gaussa]\label{th:wiener-isometry}
Dla dowolnych $f,g\in\calS_T$ zachodzą
\[
\E I_W(f)=0,\qquad
\E[I_W(f)I_W(g)]=\int_0^T f(t)g(t)\,\dd t.
\tag{1.2}\label{eq:wienscalar}
\]
Ponadto $I_W(f)$ ma rozkład normalny
$\Normal(0,\int_0^Tf(t)^2\,\dd t)$,
przy czym wariancja zero oznacza rozkład skupiony w zerze.
\end{twierdzenie}
\begin{proof}
Z poprzedniego lematu możemy przyjąć wspólny podział,
tak że $f=\sum a_k\ind{(t_k,t_{k+1}]}$ i
$g=\sum b_k\ind{(t_k,t_{k+1}]}$. Oznaczmy
$\Delta_kW=W_{t_{k+1}}-W_{t_k}$. Ponieważ przyrosty Wienera
mają średnią zero, $\E I_W(f)=0$. Dla $j\ne k$
niezależność i zerowa średnia przyrostów dają
$\E[\Delta_jW\Delta_kW]=0$; dla $j=k$ mamy
$\E(\Delta_kW)^2=t_{k+1}-t_k$. Zatem
\[
\E[I_W(f)I_W(g)]
=\sum_{j,k}a_jb_k\E[\Delta_jW\Delta_kW]
=\sum_k a_kb_k(t_{k+1}-t_k)
=\int_0^Tfg\,\dd t.
\]
Na koniec przyrosty $\Delta_kW$ są niezależnymi zmiennymi
normalnymi, a ich kombinacja liniowa jest normalna.
Obliczona wariancja określa dokładnie jej rozkład.
\end{proof}

\begin{uwaga}[Ograniczenie do funkcji deterministycznych]
W powyższym dowodzie wolno było wyciągnąć współczynniki
$a_k$ i $b_k$ przed wartość oczekiwaną, ponieważ były liczbami.
Dla losowych współczynników nie jest to dopuszczalne.
W rozdziale drugim użyjemy w zamian warunkowania względem
$\calF_{t_k}$ oraz niezależności przyszłego przyrostu od tej
$\sigma$-algebry. Właśnie to prowadzi do przewidywalności.
\end{uwaga}

\section{Domknięcie przestrzeni: definicja dla funkcji kwadratowalnych}
\begin{lemat}[Gęstość funkcji schodkowych]
Przestrzeń $\calS_T$ jest gęsta w $L^2([0,T])$.
\end{lemat}
\begin{proof}
Wskaźniki przedziałów półotwartych generują $\calB([0,T])$.
Ich skończone kombinacje liniowe są gęste w $L^2$:
po pierwsze, wskaźniki zbiorów borelowskich można
aproksymować w $L^2$ wskaźnikami skończonych sum przedziałów
(z regularności miary Lebesgue'a); po drugie,
każdą funkcję z $L^2$ przybliża się funkcjami prostymi,
najpierw przez obcięcie, a potem przez kwantyzację jej wartości.
Po wspólnym rozdrobnieniu przedziałów wszystkie aproksymacje
należą do $\calS_T$.
\end{proof}

\begin{twierdzenie}[Całka Wienera w $L^2$]\label{th:wiener-extension}
Istnieje dokładnie jedno liniowe i ciągłe odwzorowanie
$I_W:L^2([0,T])\to L^2(\Omega)$, które zgadza się z
\eqref{eq:wiener-simple} na $\calS_T$. Jest izometrią:
\[
\norm{I_W(f)}_{L^2(\Omega)}=\norm{f}_{L^2([0,T])}.
\tag{1.3}\label{eq:wiener-iso}
\]
W szczególności $I_W(f)$ jest granicą w $L^2(\Omega)$
zmiennych $I_W(f_n)$ dla dowolnego ciągu $f_n\in\calS_T$
zbiegającego do $f$ w $L^2([0,T])$. Definicja granicy
nie zależy od wyboru ciągu aproksymującego.
\end{twierdzenie}
\begin{proof}
Z izometrii funkcji schodkowych dostajemy dla $m,n$
$\norm{I_W(f_n)-I_W(f_m)}_2=\norm{f_n-f_m}_2$.
Skoro $f_n\to f$ w $L^2([0,T])$, ciąg $(f_n)$ jest Cauchy'ego,
a więc $(I_W(f_n))$ jest Cauchy'ego w zupełnej przestrzeni
$L^2(\Omega)$. Ma zatem granicę $Y$. Jeżeli $g_n\to f$
jest innym ciągiem schodkowym, to
$\norm{I_W(f_n)-I_W(g_n)}_2=\norm{f_n-g_n}_2\to0$,
więc oba ciągi mają tę samą granicę. Kładziemy $I_W(f)=Y$.
Liniowość wynika z liniowości dla funkcji schodkowych i
ciągłości dodawania w $L^2$. Przechodząc do granicy w izometrii
otrzymujemy \eqref{eq:wiener-iso}; dowodzi to także ciągłości.
Jeżeli inne ciągłe liniowe odwzorowanie zgadzałoby się na gęstej
podprzestrzeni $\calS_T$, to przez ciągłość zgadzałoby się
wszędzie.
\end{proof}

\begin{twierdzenie}[Gaussowskość i kowariancja całki]\label{th:wiener-gaussian}
Dla dowolnych $f_1,\ldots,f_m\in L^2([0,T])$ wektor
$(I_W(f_1),\ldots,I_W(f_m))$ jest centrowanym wektorem gaussowskim
z macierzą kowariancji
$\big(\int_0^T f_i(t)f_j(t)\,\dd t\big)_{i,j=1}^m$.
W szczególności
\[
\E e^{iuI_W(f)}=\exp\left(-\frac{u^2}{2}\norm{f}_2^2\right),
\quad u\in\R.
\tag{1.4}
\]
Jeżeli $\int_0^Tfg\,\dd t=0$, to $I_W(f)$ i $I_W(g)$
są niezależne.
\end{twierdzenie}
\begin{proof}
Dla każdej kombinacji liniowej $h=\sum_{j=1}^ma_jf_j$
zmienna $I_W(h)$ jest granicą w $L^2$, a więc według
rozkładu, zmiennych normalnych $I_W(h_n)$ o wariancjach
dążących do $\norm{h}_2^2$. Zbieżność funkcji
charakterystycznych daje $I_W(h)\sim\Normal(0,\norm{h}_2^2)$.
Każda kombinacja liniowa wektora całek jest więc normalna,
co stanowi definicję wspólnej gaussowskości. Wzór na kowariancję
wynika z polaryzacji izometrii:
$\E I_W(f)I_W(g)=\int_0^Tfg\,\dd t$.
Dla wektora gaussowskiego nieskorelowanie dwóch współrzędnych
implikuje niezależność, co daje ostatnie stwierdzenie.
\end{proof}

\begin{wniosek}[Całkowanie po rozłącznych przedziałach]
Jeżeli $f,g\in L^2([0,T])$ mają rozłączne nośniki z dokładnością
do zbioru miary Lebesgue'a zero, to $I_W(f)$ i $I_W(g)$
są niezależne. Dla $0\le a<b\le T$ mamy
$I_W(\ind{(a,b]})=W_b-W_a$.
\end{wniosek}

\begin{przyklad}[Od funkcji stałej do funkcji wykładniczej]
Dla $f\equiv1$ mamy $I_W(f)=W_T\sim\Normal(0,T)$.
Dla $f=\ind{(a,b]}$ otrzymujemy przyrost Wienera.
Dla $f(t)=e^{-\lambda t}$, $\lambda>0$, wynik jest
zmienną normalną o średniej zero i wariancji
\[
\int_0^Te^{-2\lambda t}\,\dd t
=\frac{1-e^{-2\lambda T}}{2\lambda}.
\]
Łatwo w ten sposób znaleźć rozkład całki mimo braku
klasycznego całkowania wzdłuż trajektorii.
\end{przyklad}

\begin{zadanie}[Rozkład całki i niezależność]
Obliczyć prawo losowej zmiennej
$X=\int_0^T(2t-T)\,\dd W_t$ oraz ustalić,
czy jest ona niezależna od $W_T$.

\emph{Rozwiązanie.}
Z twierdzenia \ref{th:wiener-gaussian} zmienna $X$ ma rozkład
normalny $\Normal(0,\sigma^2)$, gdzie
$\sigma^2=\int_0^T(2t-T)^2\,\dd t=T^3/3$.
Ponadto
$\Cov(X,W_T)=\int_0^T(2t-T)\cdot1\,\dd t=0$.
Para $(X,W_T)$ jest wspólnie gaussowska, zatem jej
współrzędne są niezależne. Samo zerowanie się kowariancji
bez wspólnej gaussowskości nie wystarczyłoby do tego wniosku.
\end{zadanie}


\begin{zadanie}[Osobliwość funkcji podcałkowej]
Ustalić, dla jakich $\alpha\in\R$ istnieje całka Wienera
$X_\alpha=\int_0^Tt^{-\alpha}\,\dd W_t$, a dla jakich
nie może powstać jako granica całek obciętych
$\int_\varepsilon^Tt^{-\alpha}\,\dd W_t$
w prawdopodobieństwie, gdy $\varepsilon\downarrow0$.

\emph{Rozwiązanie.} Funkcja $t^{-\alpha}$ należy do
$L^2([0,T])$ wtedy i tylko wtedy, gdy
$\int_0^Tt^{-2\alpha}\,\dd t<\infty$, czyli
$2\alpha<1$. Dla $\alpha<1/2$ całka istnieje,
a jej rozkład jest normalny ze średnią zero i wariancją
$T^{1-2\alpha}/(1-2\alpha)$.

Jeżeli $\alpha\ge1/2$, to dla każdego $\varepsilon>0$
zmienna $X_\varepsilon:=\int_\varepsilon^T t^{-\alpha}\,\dd W_t$
jest normalna z wariancją
$v_\varepsilon=\int_\varepsilon^Tt^{-2\alpha}\,\dd t$,
która dąży do nieskończoności. Ponieważ
$X_\varepsilon/\sqrt{v_\varepsilon}\sim\Normal(0,1)$,
dla każdego ustalonego $R>0$ zachodzi
$\Prob(|X_\varepsilon|\le R)
=\Prob(|Z|\le R/\sqrt{v_\varepsilon})\to0$.
Rodzina rozkładów $(X_\varepsilon)$ nie jest ciasna,
a więc nie może zbiegać w prawdopodobieństwie do
skończonej zmiennej losowej. W ten sposób warunek
$f\in L^2$ otrzymuje także konkretne znaczenie:
zapobiega nieograniczonemu wzrostowi wariancji.
\end{zadanie}

\begin{przyklad}[Ziarno wykładnicze i początek procesu Ornsteina--Uhlenbecka]
Dla $\lambda>0$ i $t\in[0,T]$ rozważmy
$X_t:=\int_0^te^{-\lambda(t-s)}\,\dd W_s$.
Dla każdej ustalonej chwili $t$ funkcja
$s\mapsto e^{-\lambda(t-s)}\ind{(0,t]}(s)$
jest deterministyczna i należy do $L^2([0,T])$.
Stąd $X_t$ jest centrowaną zmienną normalną, a
\[
\Var(X_t)=\int_0^te^{-2\lambda(t-s)}\,\dd s
=\frac{1-e^{-2\lambda t}}{2\lambda}.
\]
Dla $0\le s\le t$ izometria spolaryzowana daje
\[
\Cov(X_s,X_t)
=\int_0^s e^{-\lambda(s-u)}e^{-\lambda(t-u)}\,\dd u
=e^{-\lambda(t-s)}\frac{1-e^{-2\lambda s}}{2\lambda}.
\]
Otrzymaliśmy bez rozwiązywania równania stochastycznego
konkretny proces gaussowski z pamięcią wykładniczą.
W dalszych rozdziałach pokażemy, że jest to rozwiązanie
liniowego równania stochastycznego z dryfem
$-\lambda X_t$.
\end{przyklad}

\section{Proces całkowy dla deterministycznej funkcji}
\begin{definicja}[Proces całek Wienera]
Dla $f\in L^2([0,T])$ i $t\in[0,T]$ kładziemy
$M_t^f:=I_W(f\ind{(0,t]})$.
Zmienna $M_t^f$ jest dobrze określona w $L^2(\Omega)$.
\end{definicja}
\begin{twierdzenie}[Martyngał, kowariancja i przyrosty]
Proces $(M_t^f)_{0\le t\le T}$ jest martyngałem względem
$(\calF_t)$ i ma ciągłą modyfikację. Dla $0\le s\le t\le T$
zachodzą
\[
\E[M_t^f\mid\calF_s]=M_s^f,\qquad
\E[M_s^fM_t^f]=\int_0^sf(u)^2\,\dd u.
\tag{1.5}
\]
Przyrost $M_t^f-M_s^f$ jest niezależny od $\calF_s$
i ma rozkład $\Normal(0,\int_s^tf(u)^2\,\dd u)$.
\end{twierdzenie}
\begin{proof}
Najpierw załóżmy, że $f$ jest schodkowa. Przyrost
$M_t^f-M_s^f$ jest deterministyczną kombinacją przyrostów
procesu Wienera na $(s,t]$. Jest więc niezależny od $\calF_s$
i ma średnią zero. Zatem własność martyngałowa zachodzi.
Dla ogólnego $f$ wybierzmy $f_n\in\calS_T$ z
$\norm{f_n-f}_2\to0$. Z izometrii wynika
$M_t^{f_n}\to M_t^f$ w $L^2$ dla każdego $t$.
Warunkowa wartość oczekiwana jest kontrakcją w $L^1$,
więc przechodząc do granicy otrzymujemy
$\E[M_t^f\mid\calF_s]=M_s^f$. Jednocześnie rozkład
i niezależność przyrostu od $\calF_s$ wynikają z granicy
w $L^2$ oraz faktu, że zmienne niezależne od danej
$\sigma$-algebry zachowują tę własność przy zbieżności
według prawdopodobieństwa (wystarczy sprawdzić równość
oczekiwań $\E[g(Y_n)\ind A]=\E g(Y_n)\Prob(A)$ dla
$A\in\calF_s$, a następnie przejść do granicy dla
granicznie ograniczonych ciągłych $g$). Wzór na kowariancję
wynika z \eqref{eq:wienscalar} po domknięciu.

Ciągłość trajektorii wymaga osobnego argumentu. Dla
schodkowej $f_n$ proces $M^{f_n}$ jest ciągły, a
nierówność maksymalna Dooba i izometria dają
\[
\E\sup_{0\le t\le T}|M_t^{f_n}-M_t^{f_m}|^2
\le4\norm{f_n-f_m}_2^2.
\]
Możemy więc wybrać podciąg, dla którego
$\sum_j\big(\E\sup_t|M_t^{f_{n_{j+1}}}-M_t^{f_{n_j}}|^2\big)^{1/2}<\infty$.
Nierówność Cauchy'ego--Schwarza i twierdzenie o monotonicznej
zbieżności pokazują, że suma supremów różnic jest skończona
prawie na pewno. Wobec tego podciąg zbiega jednostajnie
prawie na pewno do procesu o ciągłych trajektoriach.
Dla każdego ustalonego $t$ granica zgadza się z $M_t^f$
w $L^2$, a więc jest jego modyfikacją.
\end{proof}

\begin{uwaga}[Zmiana czasu]
Funkcja $A_f(t):=\int_0^tf(u)^2\,\dd u$ jest ciągła
i niemalejąca. Rozkłady skończeniewymiarowe procesu $M^f$
są takie same jak procesu $B_{A_f(t)}$, gdzie $B$
jest standardowym ruchem Browna na odpowiednio rozszerzonej
przestrzeni czasu. Jest to równość \emph{praw procesów},
a nie stwierdzenie $M_t^f=B_{A_f(t)}$ dla ustalonego
wcześniej, niezależnego ruchu Browna.
\end{uwaga}

\section{Całkowanie przez części i granica sum}
\begin{twierdzenie}[Deterministyczne całkowanie przez części]
Jeżeli $f:[0,T]\to\R$ jest funkcją absolutnie ciągłą
i $f'\in L^2([0,T])$, to
\[
\int_0^T f(t)\,\dd W_t
=f(T)W_T-\int_0^T f'(t)W_t\,\dd t
\quad\text{prawie na pewno}.
\tag{1.6}\label{eq:ibp-det}
\]
Całka po prawej jest zwykłą całką Lebesgue'a względem czasu.
\end{twierdzenie}
\begin{proof}
Dla funkcji schodkowej wyznaczonej wartościami $f(t_k)$
suma Wienera spełnia dyskretne całkowanie przez części:
\[
\sum_{k=0}^{n-1}f(t_k)\Delta_kW
=f(T)W_T-\sum_{k=0}^{n-1}W_{t_{k+1}}(f(t_{k+1})-f(t_k))
\]
jest dokładną tożsamością teleskopową, ponieważ $W_0=0$.
Wystarczy zatem zbadać granice obu stron. Funkcje $f_n(t)=f(t_k)$ dla
$t\in(t_k,t_{k+1}]$ zbiegają jednostajnie do $f$,
a więc w $L^2([0,T])$. Lewa strona zbiega zatem do
$I_W(f)$ w $L^2(\Omega)$.
Ponieważ $f$ jest absolutnie ciągła,
$\Delta_k f=\int_{t_k}^{t_{k+1}}f'(u)\,\dd u$,
a więc druga suma jest równa
$\int_0^T W_{r_n(u)}f'(u)\,\dd u$,
gdzie $r_n(u)$ jest prawym końcem przedziału zawierającego $u$.
Ciągłość trajektorii $W$ zapewnia jednostajną zbieżność
$W_{r_n(u)}\to W_u$ prawie na pewno. Ponieważ
$\int_0^T|f'(u)|\,\dd u<\infty$, suma zbiega niemal
na pewno do $\int_0^TW_uf'(u)\,\dd u$.
Granice w prawdopodobieństwie są jednoznaczne, co daje tezę.
\end{proof}

\begin{przyklad}[Ważona całka z procesu Wienera]
Dla $f(t)=t$ otrzymujemy
$\int_0^Tt\,\dd W_t=T W_T-\int_0^TW_t\,\dd t$.
Losowa zmienna po lewej ma rozkład
$\Normal(0,\int_0^Tt^2\,\dd t)=\Normal(0,T^3/3)$.
Nie mylimy jej ze zwykłą całką $\int_0^TW_t\,\dd t$.
Co ciekawe, obie zmienne mają ten sam rozkład normalny
$\Normal(0,T^3/3)$, a nawet tę samą kowariancję
$T^2/2$ z $W_T$. Nie są jednak równe:
na mocy \eqref{eq:ibp-det} ich różnica to
$\int_0^T(2t-T)\,\dd W_t$, mająca wariancję
$T^3/3>0$. To przykład pokazujący, że jednakowy
rozkład nie oznacza równości zmiennych losowych.
\end{przyklad}

\section{Miara Wienera a całka po trajektoriach}
Dotychczas całkowaliśmy funkcję czasu względem losowych
przyrostów. Istnieje jeszcze inny rodzaj całkowania:
całkowanie funkcjonałów całej ścieżki względem \emph{miary
Wienera}. Nie wolno tych pojęć utożsamiać.

\begin{definicja}[Miara Wienera i całka po trajektoriach]
Niech $C_0([0,T])$ oznacza przestrzeń funkcji ciągłych
$w:[0,T]\to\R$ o $w(0)=0$, z normą supremum i borelowskim
$\sigma$-ciałem. Niech $\mu_W$ będzie rozkładem procesu Wienera
jako zmiennej losowej o wartościach w $C_0([0,T])$:
$\mu_W(A):=\Prob(W\in A)$ dla $A\in\calB(C_0([0,T]))$.
Dla mierzalnego i całkowalnego funkcjonału
$F:C_0([0,T])\to\R$ jego \emph{całką względem miary Wienera}
nazywamy
\[
\int_{C_0([0,T])}F(w)\,\dd\mu_W(w)=\E[F(W)].
\tag{1.7}
\]
\end{definicja}

\begin{uwaga}[Różne obiekty po dwóch stronach znaku całki]
W $\int_0^Tf(t)\,\dd W_t$ zmienna $t$ przebiega przedział
czasu, a wynikiem jest zmienna losowa. W
$\int_{C_0([0,T])}F(w)\,\dd\mu_W(w)$ przebiegamy zbiór
całych funkcji $w$ i otrzymujemy liczbę. Funkcjonał $F$
może zależeć na przykład od maksimum trajektorii,
jej pola pod wykresem albo wartości w wielu chwilach.
Zapis $\dd\mu_W(w)$ jest ścisłym zapisem całki po trajektoriach
w czasie euklidesowym. Formalna całka Feynmana o oscylacyjnym
czynniku $\exp(iS/\hbar)$ wymaga innej konstrukcji.
\end{uwaga}

\begin{zadanie}[Całka funkcyjna dla pola pod trajektorią]\label{zad:area-char}
Obliczyć dla $u\in\R$ całkę
\[
\int_{C_0([0,T])}
\exp\left(iu\int_0^T w(s)\,\dd s\right)\,\dd\mu_W(w).
\]
\emph{Rozwiązanie.}
Kładziemy $A_T=\int_0^TW_s\,\dd s$.
Z \eqref{eq:ibp-det} zastosowanego do $f(s)=T-s$
wynika $A_T=I_W(T-\cdot)$, ponieważ $(T-s)'=-1$
i $f(T)=0$. Jest to zmienna normalna o średniej zero
i wariancji
$\int_0^T(T-s)^2\,\dd s=T^3/3$.
Zatem jej funkcja charakterystyczna wynosi
\[
\E e^{iuA_T}=\exp\left(-\frac{u^2T^3}{6}\right).
\]
Otrzymaliśmy całkę po nieskończeniewymiarowej przestrzeni
trajektorii, sprowadzając ją do funkcji charakterystycznej
jednej normalnej zmiennej losowej.
\end{zadanie}

\begin{przyklad}[Całka cylindryczna]
Niech $0<t_1<\cdots<t_m\le T$ i
$F(w)=\varphi(w(t_1),\ldots,w(t_m))$ dla ograniczonej
funkcji borelowskiej $\varphi:\R^m\to\R$.
Wówczas $\int F\,\dd\mu_W$ jest zwykłą
$m$-wymiarową całką względem wspólnego rozkładu
$(W_{t_1},\ldots,W_{t_m})$. Jest to rozkład gaussowski
z macierzą kowariancji $K_{ij}=\min(t_i,t_j)$.
Dla przykładów nieskończeniewymiarowych naturalną metodą
jest aproksymowanie funkcjonału takimi funkcjonałami
cylindrycznymi, a następnie uzasadnienie zbieżności całek.
\end{przyklad}

\section{Dlaczego całka Wienera nie jest punktowym funkcjonałem ciągłym?}
Izometria \eqref{eq:wiener-iso} zachodzi pomiędzy
\emph{przestrzeniami klas funkcji}. Nie mówi, że dla prawie
każdego $\omega$ odwzorowanie
$f\mapsto I_W(f)(\omega)$ jest ograniczonym funkcjonałem
liniowym na $L^2([0,T])$. Takie stwierdzenie byłoby fałszywe.

\begin{twierdzenie}[Brak jednoczesnej ścieżkowej ciągłości]
Nie istnieje zbiór $\Omega_0\subset\Omega$ o
$\Prob(\Omega_0)=1$ i rodzina ograniczonych liniowych
funkcjonałów $L_\omega:L^2([0,T])\to\R$, $\omega\in\Omega_0$,
taka że dla każdego $f\in L^2([0,T])$ równość
$L_\omega(f)=I_W(f)(\omega)$ zachodzi niemal na pewno
(w sensie wybranych wersji zmiennych $I_W(f)$).
\end{twierdzenie}
\begin{proof}
Niech $(e_n)_{n\ge1}$ będzie układem ortonormalnym zupełnym
w $L^2([0,T])$; taka baza istnieje, ponieważ jest to
separowalna przestrzeń Hilberta. Z twierdzenia
\ref{th:wiener-gaussian} wynika, że
$Z_n:=I_W(e_n)$ są niezależnymi zmiennymi o rozkładzie
$\Normal(0,1)$. Ponieważ baza jest przeliczalna, możemy
wybrać wspólny zbiór pełnego prawdopodobieństwa, na którym
$L_\omega(e_n)=Z_n(\omega)$ dla wszystkich $n$.
Jeżeli $L_\omega$ jest ograniczonym funkcjonałem liniowym,
na mocy twierdzenia Riesza istnieje $h_\omega\in L^2([0,T])$
taka, że $L_\omega(f)=\langle f,h_\omega\rangle$.
Nierówność Bessela daje wtedy
$\sum_{n\ge1}Z_n(\omega)^2\le\norm{h_\omega}_2^2<\infty$.
Tymczasem mocne prawo wielkich liczb dla niezależnych
zmiennych $Z_n^2$ zapewnia
$n^{-1}\sum_{k=1}^nZ_k^2\to\E Z_1^2=1$ prawie na pewno.
Oznacza to, że suma $\sum Z_n^2$ jest nieskończona
prawie na pewno, sprzeczność.
\end{proof}

\begin{uwaga}[Intuicja: biały szum]
Zapis heurystyczny $\dd W_t=\xi(t)\,\dd t$ sugeruje,
że całkujemy $f$ względem \emph{białego szumu} $\xi$.
Powyższe twierdzenie wyjaśnia, dlaczego $\xi$ nie może
być zwykłą losową funkcją z $L^2([0,T])$. Można jednak
zdefiniować biały szum jako losowy obiekt uogólniony,
który działa na funkcjach testowych; ścisłą teorię
dystrybucji losowych rozwiniemy w dalszej części tomu.
\end{uwaga}

\section{Podsumowanie konstrukcji i zadania}
Konstrukcja całki Wienera opiera się na trzech zasadach:
niezależnych normalnych przyrostach, izometrii w $L^2$
i zupełności tej przestrzeni. Umożliwia ona całkowanie
funkcji deterministycznych $f\in L^2$, lecz nie wystarcza
dla losowej funkcji podcałkowej zależnej od wartości
procesu w chwili $t$. Tę trudność rozwiążemy przez
warunkowanie względem informacji dostępnej do czasu $t$.

\begin{zadanie}[Niezależne całki z niepokrywających się przedziałów]
Dla $0<a<b<c<T$ wyznaczyć rozkład wektora
$(X,Y)$, gdzie $X=\int_a^b t\,\dd W_t$,
$Y=\int_b^c (t-b)\,\dd W_t$. Czy jego współrzędne
są niezależne?

\emph{Rozwiązanie.} Jest to wektor centrowany gaussowski.
Wariancje wynoszą $\Var(X)=(b^3-a^3)/3$ oraz
$\Var(Y)=(c-b)^3/3$, a kowariancja jest równa zero,
gdyż iloczyn funkcji podcałkowych znika niemal wszędzie.
Stąd $X$ i $Y$ są niezależne. Wspólny punkt $b$ nie
wpływa na wynik, ponieważ ma zerową miarę Lebesgue'a.
\end{zadanie}

